\section{Hydraulische Charakterisierung der Komponenten}
\label{sec:hydraulik}

Um den Druckverlust des Kühlkreislaufs und den benötigten Druckaufbau der Pumpe
bestimmen zu können, müssen zunächst die Druckverluste der einzelnen Bauteile
ermittelt werden.

\subsection{Messaufbau}
\label{subsec:hydraulik_messaufbau}

Die Messungen erfolgten am Fahrzeug RP25e/d. Dieses besitzt zwei getrennte
Kühlkreisläufe: einen für den Inverter und einen für die vier in Reihe
geschalteten Motoren.

% FIGURE-VORSCHLAG 19: Fließbild der beiden Kühlkreisläufe des RP25e/d
% (Inverter-Kreis; Motor-Kreis mit vier Motoren in Reihe, Radiatoren, Pumpen).

\subsection{Vorgehen}
\label{subsec:hydraulik_vorgehen}

Zunächst wurden beide Kreisläufe im Ausgangszustand vermessen. Dazu wurde
der Volumenstrom ermittelt, indem die innerhalb einer Minute geförderte
Wassermenge bestimmt wurde. Anschließend wurde jeweils eine Komponente aus dem
Kreis entfernt und die Messung wiederholt. Aus den unterschiedlichen
Durchflüssen lässt sich der hydraulische Widerstand der entfernten Komponente
bestimmen.

Betrachtet wurden folgende Komponenten: Schlauch mit \SI{10}{mm}
\TODO{Innen- oder Außendurchmesser} je \SI{1}{m} Länge, Motor, Radiator und
Inverter.

\subsection{Berechnung der Widerstandswerte}
\label{subsec:hydraulik_berechnung}

Für den turbulent durchströmten Kreis wird der Druckverlust als quadratisch
vom Volumenstrom abhängig angenommen:
\begin{equation}
    \Delta p = R \, \dot{V}^{2}.
    \label{eq:druckverlust}
\end{equation}
Bei Reihenschaltung addieren sich die Widerstände der Komponenten zu
$R_\mathrm{ges} = \sum_i R_i$. Die Pumpe arbeitet in dem Betriebspunkt, in dem
ihr Druckaufbau $\Delta p_\mathrm{P}(\dot{V})$ dem Druckverlust des Kreises
entspricht. Für den Ausgangszustand (Index 1) und den Zustand ohne Komponente
$K$ (Index 2) gilt damit
\begin{align}
    \Delta p_\mathrm{P}(\dot{V}_1) &= R_\mathrm{ges}\,\dot{V}_1^{2}, \\
    \Delta p_\mathrm{P}(\dot{V}_2) &= \left(R_\mathrm{ges} - R_K\right)\dot{V}_2^{2}.
\end{align}
Daraus folgt für den Widerstand der Komponente
\begin{equation}
    R_K = \frac{\Delta p_\mathrm{P}(\dot{V}_1)}{\dot{V}_1^{2}}
        - \frac{\Delta p_\mathrm{P}(\dot{V}_2)}{\dot{V}_2^{2}}.
    \label{eq:widerstand_komponente}
\end{equation}
Der Druckaufbau der Pumpe wird dabei aus der bekannten Pumpenkennlinie bei
den gemessenen Volumenströmen abgelesen.

\begin{table}[htbp]
    \centering
    \caption{Ermittelte hydraulische Widerstände der Komponenten}
    \label{tab:hydraulik_widerstaende}
    \begin{tabular}{l S[table-format=2.2] l}
        \toprule
        Komponente & {$R$ in \si{...}} & Bemerkung \\
        \midrule
        Schlauch (\SI{10}{mm}, je \SI{1}{m}) & {\TODO{}} & \\
        Motor & {\TODO{}} & \\
        Radiator & {\TODO{}} & \\
        Inverter & {\TODO{}} & \\
        \bottomrule
    \end{tabular}
\end{table}

% FIGURE-VORSCHLAG 20: Pumpenkennlinie mit den gemessenen Betriebspunkten
% (Ausgangszustand und je ein Punkt pro entfernter Komponente) und den
% zugehörigen Anlagenkennlinien.

\subsection{Modellierung und Layoutvergleich}
\label{subsec:hydraulik_modell}

Anschließend wurde der Kühlkreislauf in MATLAB nachgebildet. Die berechneten
Werte wurden mit einem Korrekturfaktor an die Messwerte angeglichen. Der
Korrekturfaktor beträgt \TODO{} und ergibt sich aus \TODO{Definition}.
Danach wurden verschiedene Layouts in MATLAB aufgebaut, deren Durchflüsse
und mögliche Wärmeabfuhr berechnet und verglichen.

% FIGURE-VORSCHLAG 21: Vergleich Messung und Modell (Volumenstrom je Konfiguration).
% FIGURE-VORSCHLAG 22: Balkendiagramm Durchfluss/Wärmeabfuhr je Layout.